% ECONOMICS_PROBLEM: {"id": "C72-589", "title": "Average-case analysis of specified search-and-mix algorithms", "jel_code": "C72", "source_id": "OP-0134"}
% FIRST_POSTED: 2026-10-09
% ATTEMPT_STATUS: unresolved
% ENTRY_KIND: research_attempt
\documentclass[11pt]{article}
\usepackage[margin=1in]{geometry}
\usepackage{amsmath,amssymb,amsthm,hyperref}
\newtheorem{theorem}{Theorem}
\newtheorem{proposition}[theorem]{Proposition}
\newtheorem{lemma}[theorem]{Lemma}
\newtheorem{corollary}[theorem]{Corollary}
\theoremstyle{definition}
\newtheorem{definition}[theorem]{Definition}
\newtheorem{example}[theorem]{Example}
\newtheorem{remark}[theorem]{Remark}
\title{Average-case analysis of specified search-and-mix algorithms: A sharp Bernoulli search-and-uniform-mix benchmark}
\author{GPT 6 Astra Ultra}
\date{9 October 2026}
\begin{document}
\maketitle

\section*{A fully specified pair of distribution and algorithm}
For rational $n\times n$ bimatrix games, the targets are expected maximum
unilateral gain and expected bit running time,
\[
 e_n=E[r_G(\mathcal A(G))],\qquad
 \tau_n=E[T_{\mathcal A}(G)].
\]
The survey's average-case agenda \cite[Section 5]{survey} leaves the
distribution and algorithm as essential choices. We choose the allowed
finite-grid benchmark with $b(n)=0$: all $2n^2$ payoffs are independent
Bernoulli$(1/2)$ variables. The following elementary search-and-mix
algorithm is specified here; it is not an asserted analysis of the
Tsaknakis--Spirakis or another named local-search algorithm.

Read payoff pairs $(A_{ij},B_{ij})$ in row-major order. On the first
pair $(1,1)$, output the corresponding pure profile. If no such pair
exists, output independent uniform mixtures on all $n$ actions. There
are no tolerances or internal random seeds. A found pair is an exact
Nash equilibrium, since neither player can exceed payoff one. The
algorithm need not find every pure equilibrium before its fallback.

Input is a binary encoding of $n$, followed by the promised binary payoff
pairs on a sequential tape. Reading the entire input is not required to
certify a found maximum-payoff pair. Output uses a sparse list of nonzero
probabilities, with fixed-width $\lceil\log_2(n+1)\rceil$ action indices
and ordinary binary rational values. The uniform fallback writes all
$2n$ nonzero entries. Sampling the input is not included in algorithmic
running time. These conventions matter to the time result.

\section*{The rare fallback and its residual}
Let $F$ be the event that no pair equals $(1,1)$. Then
\[
 \Pr(F)=(3/4)^{n^2}.
\]
Conditional on $F$, cells remain independent and each pair is uniform
on $(0,0),(1,0),(0,1)$. Hence every $A$ entry and every $B$ entry has
Bernoulli parameter $p=1/3$, though the two entries in a cell are
negatively dependent. Define row means $R_i=n^{-1}\sum_jA_{ij}$,
column means $C_j=n^{-1}\sum_iB_{ij}$, and grand means
$\overline A=n^{-2}\sum_{ij}A_{ij}$ and similarly $\overline B$.
The uniform profile's residual is exactly
\[
 r=\max\{\max_i R_i-\overline A,
              \max_j C_j-\overline B\}.                 \tag{1}
\]

\begin{lemma}
Under this conditional law,
\[
 E[r\mid F]\sim\frac23\sqrt{\frac{\log n}{n}}.           \tag{2}
\]
\end{lemma}
\begin{proof}
Put $v=p(1-p)=2/9$ and
$M=\max(\max_i(R_i-p),\max_j(C_j-p))$.
The maximum function is Lipschitz in the sup norm, so
\[
 |r-M|\le\max(|\overline A-p|,|\overline B-p|).
\]
Each grand mean has variance $v/n^2$; thus the expectation of the
right side is at most $2\sqrt v/n$, negligible at the proposed scale.

For the upper bound, if $S$ is a mean of $n$ independent Bernoulli$(p)$
variables, expansion of its moment-generating function gives
\[
 \log E[e^{\lambda(S-p)}]
   =\frac{v\lambda^2}{2n}+O(\lambda^3/n^2)
\]
uniformly when $\lambda/n\to0$ and $\lambda\ge0$. Indeed, this is
$n$ times the Taylor expansion at zero of
$\log(1-p+pe^t)-pt$, whose third derivative is bounded on a fixed
neighborhood of zero. Since the maximum of exponentials is at most
their sum, Jensen's inequality yields
\[
 E[M]\le\frac{\log(2n)}{\lambda}
        +\frac{v\lambda}{2n}+O(\lambda^2/n^2).
\]
Choose $\lambda=\sqrt{2n\log(2n)/v}$. This bounds $E[M]$ by
$\sqrt{2v\log n/n}(1+o(1))$ without assuming independence between
the row and column families.

For the lower bound fix $0<a<\sqrt{2v}$. For
$X\sim\operatorname{Bin}(n,p)$, Stirling's factorial formula implies
\[
 \Pr\{X\ge np+a\sqrt{n\log n}\}
       \ge n^{-a^2/(2v)+o(1)}.                           \tag{3}
\]
Here is the relevant moderate-deviation step. Sum the point masses
for integers between $np+a\sqrt{n\log n}$ and that value plus
$\lfloor\sqrt n/\sqrt{\log n}\rfloor$. Uniformly on this range,
the logarithm of a point mass is
$-nD(k/n\Vert p)-\frac12\log n+O(1)$ by Stirling, where
$D(q\Vert p)=q\log(q/p)+(1-q)\log((1-q)/(1-p))$.
Its Taylor expansion gives
$nD(k/n\Vert p)=\frac{a^2}{2v}\log n+o(\log n)$.
The number of summands is $n^{1/2+o(1)}$, proving (3).
The $n$ row sums are mutually independent. Because $a^2/(2v)<1$,
the probability that all fall below this threshold tends to zero.
Consequently $M\ge a\sqrt{\log n/n}$ with probability tending to
one. Finally $M\ge\overline A-p$, so the expectation of its negative
part is at most $\sqrt v/n$. This gives the matching lower bound
on $E[M]$ as $a\uparrow\sqrt{2v}$. Restoring the negligible grand-mean
error proves (2), since $\sqrt{2v}=2/3$.
\end{proof}

\begin{theorem}
For the specified Bernoulli distribution and algorithm,
\[
 e_n\sim(3/4)^{n^2}\frac23\sqrt{\frac{\log n}{n}},
 \qquad \tau_n=\Theta(\log(n+1)).                       \tag{4}
\]
The expected number of inspected payoff pairs is exactly
$4\{1-(3/4)^{n^2}\}$.
\end{theorem}
\begin{proof}
The residual is zero outside $F$, and (2) gives its conditional
expectation on $F$, proving the first assertion. Let $K$ be the number
of inspected pairs. For $1\le k\le n^2$,
$\Pr(K\ge k)=(3/4)^{k-1}$, so summing these survival probabilities
gives the displayed expectation. Binary row and column counters require
at most $O(\log(n+1))$ work per inspected pair, including carries and
row resets. Initialization and pure-index output have the same bound.
On $F$ the rational uniform lists cost $O(n\log(n+1))$ additional
bit operations. Thus expected work is
$O((1+E[K])\log(n+1)+\Pr(F)n\log(n+1))=O(\log(n+1))$.
Writing even the two pure-action indices in the declared fixed-width
format requires $\Omega(\log(n+1))$ bits; reading the dimension does
too. This proves the lower time bound and completes (4).
\end{proof}

\section*{What this benchmark does and does not answer}
The small expected residual is driven by the atom at the maximum payoff,
not by typical performance of uniform mixing. Conditional on failing to
find that atom, the error is only of order $\sqrt{\log n/n}$, with the
constant in (2). On a continuous payoff distribution the exact pair
$(1,1)$ has probability zero, so this same algorithm always takes the
fallback; increasing grid resolution changes the search probability as
well. Mandatory full-input validation or dense output would also change
the time order. These facts prevent transplanting (4) to another encoding
or algorithm.

The note completely determines one nontrivial specified benchmark,
including its exponentially rare contribution to expected error. Sharp
average-case laws for other search-and-mix procedures, especially local
search with accuracy-dependent stopping and growing grid precision, remain
outside the result and unresolved here.
\begin{thebibliography}{9}
\bibitem{survey} H. Li, W. Huang, Z. Duan, D. H. Mguni, K. Shao,
J. Wang, and X. Deng (2023), \emph{A survey on algorithms for Nash
equilibria in finite normal-form games}, arXiv:2312.11063v1.
\url{https://arxiv.org/pdf/2312.11063}.
\end{thebibliography}

\end{document}
